\chapter{HANA Schema Extraction Pipeline}
\label{chap:extraction}\label{sim:chap:extraction}

\section{Pipeline at a Glance}

\mermaidfig[\textwidth]{diagrams/extraction-pipeline}%
  {The five extraction stages; each emits a well-typed payload to the next.}%
  {fig:extraction-pipeline}

\section{Stage 1 --- MCP Discovery}

The extractor connects to the AI Core PAL MCP server and discovers available tools:

\begin{lstlisting}[language=jsonx,caption={MCP connection (hana\_schema\_extractor.py, lines 45--60)}]
async def connect(self) -> None:
    """Establish MCP connection to AI Core PAL."""
    # Handle both HTTP->WS and HTTPS->WSS
    transport_url = self.config.mcp_url
    transport_url = transport_url.replace("https://", "wss://")
    transport_url = transport_url.replace("http://", "ws://")
    transport_url = f"{transport_url}/mcp/ws"
    
    self._transport = WebSocketClientTransport(transport_url)
    self._session = ClientSession(self._transport)
    
    await self._session.initialize()
    
    # Discover available tools
    tools_result = await self._session.list_tools()
    self._available_tools = {t.name for t in tools_result.tools}
\end{lstlisting}

\textbf{Precondition.} The MCP server must expose at least one of: \texttt{execute\_sql},
\texttt{hana\_tables}, or \texttt{list\_tables}.

\section{Stage 2 --- Table Enumeration}

Tables are discovered via the \texttt{hana\_tables} MCP tool or by querying
\texttt{SYS.TABLES}:

\begin{lstlisting}[language=jsonx,caption={Table discovery (hana\_schema\_extractor.py, lines 85--105)}]
async def _discover_tables(self, schema: str) -> list[str]:
    """Discover tables in a schema."""
    if "hana_tables" in self._available_tools:
        result = await self._session.call_tool(
            "hana_tables",
            {"schema": schema}
        )
        return [t["table_name"] for t in result.content]
    else:
        # Fallback to direct SQL
        result = await self._execute_sql(
            'SELECT TABLE_NAME FROM SYS.TABLES '
            'WHERE SCHEMA_NAME = ?',
            [schema]
        )
        return [row["TABLE_NAME"] for row in result]
\end{lstlisting}

\section{Stage 3 --- Column Extraction}

For each table, column metadata is extracted from \texttt{SYS.TABLE\_COLUMNS}:

\begin{lstlisting}[language=jsonx,caption={Column extraction SQL}]
SELECT 
    COLUMN_NAME,
    DATA_TYPE_NAME,
    LENGTH,
    IS_NULLABLE
FROM SYS.TABLE_COLUMNS
WHERE SCHEMA_NAME = ? AND TABLE_NAME = ?
ORDER BY POSITION
\end{lstlisting}

\section{Stage 4 --- Registry Population}

Extracted metadata is normalized and stored in the \texttt{SchemaRegistry}:

\begin{lstlisting}[language=jsonx,caption={Registry population (hana\_schema\_extractor.py, lines 140--165)}]
async def extract_table_schema(
    self, schema: str, table: str
) -> TableSchema:
    """Extract full schema for a single table."""
    columns_data = await self._get_table_columns(schema, table)
    
    columns = [
        Column(
            name=col["COLUMN_NAME"],
            data_type=col["DATA_TYPE_NAME"],
            length=col.get("LENGTH"),
            nullable=col.get("IS_NULLABLE", "TRUE") == "TRUE",
        )
        for col in columns_data
    ]
    
    return TableSchema(
        schema_name=schema,
        table_name=table,
        columns=columns,
    )
\end{lstlisting}

\section{Stage 5 --- JSON Export}

The registry exports to JSON for caching and debugging:

\begin{lstlisting}[language=jsonx,caption={JSON export format}]
{
  "schemas": {
    "PAL_STORE": {
      "tables": {
        "ESG_METRICS": {
          "columns": [
            {"name": "METRIC_ID", "type": "INTEGER", ...},
            {"name": "SCOPE", "type": "NVARCHAR", ...}
          ]
        }
      }
    }
  },
  "extracted_at": "2026-04-18T08:00:00Z"
}
\end{lstlisting}

% =============================================================================
% CHAPTER 4: TAXONOMY GENERATION ENGINE
% =============================================================================
\newpage